\section{Method}

\subsection{Overview}

We study on-policy distillation for visual reasoning by explicitly separating visual perception from downstream reasoning.
Our central idea is to evaluate the same frozen Teacher under two different information pathways while keeping the Student-generated reasoning state fixed.

The Student always operates under the standard inference condition.
Given an image $I$ and a textual query $Q$, it generates an on-policy response
\[
\tau_S=(y_1,\ldots,y_T)
\sim
p_\theta(\cdot\mid I,Q).
\]
At token position $t$, we denote the Student-generated prefix as
\[
h_t=(Q,y_{<t}).
\]

For each Student-visited state $h_t$, we query the same frozen Teacher from two views.
In the first view, the Teacher receives the raw image and must jointly perform visual perception and reasoning.
In the second view, the raw image is removed and replaced by ground-truth structured visual evidence, such as object identities and bounding boxes.
This second view removes the need to recover the relevant scene geometry from pixels and allows the Teacher to reason directly from an externally resolved perceptual representation.

We therefore obtain
\[
p_t^{\mathrm{vis}}
=
p_\phi(\cdot\mid I,Q,y_{<t}),
\]
and
\[
p_t^{\mathrm{str}}
=
p_\phi(\cdot\mid E^*,Q,y_{<t}),
\]
where $E^*$ denotes the training-only structured evidence.

The Teacher parameters, query, Student-generated prefix, and prediction position are identical across the two evaluations.
Only the information pathway differs.
This construction lets us compare reasoning under endogenous visual perception with reasoning under externally resolved structured evidence.

\subsection{Decoupling Perception and Reasoning}

Standard visual reasoning couples two processes inside the same model:
\[
I
\rightarrow
\text{perception}
\rightarrow
\text{reasoning}.
\]

Errors in the final prediction may therefore arise from either stage.
A model may fail because it extracts an incorrect perceptual state from the image, or because it reasons incorrectly even when the relevant perceptual information is available.

To separate these effects, we introduce a structured-evidence view in which the perceptual state is externally specified:
\[
E^*
\rightarrow
\text{reasoning}.
\]

In our setting,
\[
E^*
=
\{(c_i,b_i)\}_{i=1}^{N},
\]
where $c_i$ denotes an object identity and $b_i$ its ground-truth bounding box.

The structured-evidence branch is not treated as a different Teacher.
Instead, it represents the same reasoner operating after the perceptual burden has been removed.
Comparing the two views therefore provides a controlled way to study how the Teacher's reasoning behavior changes when visual grounding is externally resolved.

We refer to this comparison as a \emph{perception-reasoning decoupling view}.

\subsection{Reasoning Disagreement Across Information Views}

The two Teacher views can produce different next-token distributions at the same Student state.
We use this difference as a signal of how sensitive the current reasoning decision is to the perceptual pathway.

Formally, we define
\[
d_t
=
D
\left(
p_t^{\mathrm{str}},
p_t^{\mathrm{vis}}
\right),
\]
where $D(\cdot,\cdot)$ denotes a distributional discrepancy measure.

When $d_t$ is small, the Teacher reaches similar reasoning preferences from the raw image and from the externally resolved structured evidence.
When $d_t$ is large, the Teacher's reasoning changes substantially once the perceptual state is replaced by oracle information.

Importantly, this discrepancy is not interpreted as a direct measurement of the Teacher's internal perception error.
The structured representation $E^*$ is an external surrogate rather than the Teacher's latent perceptual state.
Instead, $d_t$ indicates that the current reasoning decision is sensitive to whether task-relevant perception is resolved internally or supplied externally.

\subsection{Decoupled On-Policy Distillation}

The Student trajectory remains fully on-policy throughout training.
The Student receives only $(I,Q)$ and visits the same states it would encounter at inference time.
The structured evidence is used only to construct training-time Teacher supervision.

At each Student-visited state, the raw-image Teacher view provides the Teacher's behavior under the standard visual pathway, while the structured-evidence view provides a reasoning reference under externally resolved perception.

We use the relationship between these two distributions to determine how privileged reasoning supervision should be applied.
At a high level, the method follows two principles.

First, disagreement across the two Teacher views identifies reasoning states that are sensitive to perceptual grounding.
Second, the structured-evidence view provides a target for how the same reasoner behaves once the relevant perceptual information has been resolved.

The resulting training objective can be written abstractly as
\[
\mathcal{L}
=
\sum_{t=1}^{T}
\mathcal{W}
\left(
p_t^{\mathrm{vis}},
p_t^{\mathrm{str}}
\right)
\,
\mathcal{D}
\left(
\mathcal{T}
\left(
p_t^{\mathrm{vis}},
p_t^{\mathrm{str}}
\right),
p_t^S
\right),
\]
where $\mathcal{W}$ determines the importance of each Student-visited state, $\mathcal{T}$ defines the Teacher supervision derived from the two information views, and $\mathcal{D}$ is the token-level distillation loss.

This formulation intentionally leaves the weighting and target-construction mechanism general.
The key methodological contribution is the decoupled Teacher evaluation itself: the same reasoner is queried under raw visual perception and under externally resolved structured perception at the same on-policy Student state.

\subsection{Scope and Assumptions}

The structured-evidence view provides a meaningful reasoning-only condition only when $E^*$ contains sufficient information for the task.

We therefore focus on visual reasoning problems whose answer can be derived from object identity and geometric structure, including relative position, spatial ordering, containment, overlap, and relative size.

For these tasks, bounding boxes provide a compact representation of the task-relevant perceptual state.
In contrast, questions involving color, texture, fine-grained appearance, or other information not represented by bounding boxes cannot be fully resolved from $E^*$ alone.

Our current formulation is therefore designed for geometry-sufficient visual reasoning.
Extending the framework to broader multimodal reasoning would require richer structured evidence, such as object attributes, relations, or other task-specific perceptual annotations.